
\section{Soundness and Completeness}\label{ss:sound}
We now prove the procedure is sound, i.e., that if the tableau succeeds on $\varphi$, then $\varphi$ is \SAT, and complete for non-punctual TCTL, i.e., that if $\varphi$ is \SAT, then the tableau succeeds on $\varphi$. These results are supported by the following lemma.
\begin{lemma}\label{lem:zones}
Let $N_0,\dots,N_k$ be a path of $\TabGraph$ with non-empty zones.
\begin{enumerate*}[label=\emph{(\roman*)}, ref={(\roman*)}, itemjoin={{\ }}, itemjoin*={{\ }}]
\item\label{zoneitem:feasible} For every $\nu_k\in Z_k$ there are $\nu_i\in Z_i$ and delays realising
every step, positive at every \ruleref{Delay} step.
\item\label{zoneitem:stability} If a delay crosses no integer hyperplane $x=b$ with $b\le M$, every
subformula of $\varphi$ has constant truth on the open segment it sweeps.
\item\label{zoneitem:nonzeno} A cycle of $\TabGraph$ has a non-Zeno unfolding iff it is
time-progressing, and then one whose iterations elapse a constant positive amount.
\end{enumerate*}
\end{lemma}
\begin{proof}
\ref{zoneitem:feasible} Every zone operation a rule performs (intersection with a guard, reset, time-successor,
$\cong_M$-canonicalisation) admits back-propagation, non-empty precisely because the successor
valuation lies in its image~\cite{Bouyer05TimedAutomata}; thread backwards from $\nu_k$, the delay of
a \ruleref{Delay} step being positive since $\Dopen Z$ and $\Dbd Z$ lie strictly ahead of $Z$.
\ref{zoneitem:stability} No clock crosses an integer at or below $M$, so no integer part changes and no
fractional ordering wraps, while clocks above $M$ are identified by extrapolation; the segment lies in
one clock region, on which TCTL truth is
constant~\cite{AlurCourcoubetisDill1990,Bouyer05TimedAutomata}, and its endpoints are nodes of their
own by the integer split of \ruleref{Delay}.
\ref{zoneitem:nonzeno} By contraposition, a cycle elapsing no time takes zero time over any number of
traversals; and if some clock is neither reset on it nor extrapolated above $M$, it keeps a bound
$c\le M$, grows by exactly the delay elapsed, and caps the total at $c$. Conversely \ref{zoneitem:feasible} yields
a traversal elapsing some $\delta>0$, and the cycle returning to its first node, no rule distinguishes
two visits and no clock depletes a budget, so $\delta$ recurs at every iteration.
\end{proof}

\subsection{Soundness}\label{subsec:soundness}


In order to prove soundness, we proceed by constructing a witness automaton $\mathcal A_\varphi$ from a successful marking $\mathcal M^{*}$ of the tableau graph $\TabGraph$. We then show that this automaton satisfies the formula $\varphi$.

The construction is straightforward. The saturated nodes of $\mathcal M^{*}$ supply the locations
$S$ of $\mathcal A_\varphi$: a saturated $N$ carries $L(N)=\Gamma_{\mathsf{prop}}$ and
$\mathsf{Inv}(N)$, and $N_0$ is the initial location. Every \ruleref{Delay}- or
\ruleref{EX}-child of $N$ that $\mathcal M^{*}$ retains supplies an outgoing edge of $E$, leading
to the next node and carrying the reset clocks, i.e., $x_\mu$ at a \ruleref{launch} and at a \ruleref{Ustopcarry}. Those rules
elapse no time and leave $L$ and $\mathsf{Inv}$ untouched, so they contribute no edge of their own.

Each location carries a schedule, which follows the finite path of $\mathcal M^{*}$ discharging one of the at most $n+1$ obligations pending there, resolving the \ORn-nodes along that path and turning round-robin to the next obligation once the current one is discharged. Its memory is the pair (obligation, position on that path), so $S$ stays finite.

%
%\ref{M1} and \ref{M2} hold of a path, however, and a
%location must commit to one, so it carries in addition a \emph{schedule}. At most $n+1$
%obligations are pending at $N$, the divergence of \ref{M1} and the until pursuits active there, and
%each is discharged by a finite path of $\mathcal M^{*}$; the schedule follows one such path,
%resolving the \ORn-nodes along it, and turns round-robin to the next obligation once the current one
%is discharged. Its memory is the pair (obligation, position on that path), so the locations stay
%finite.
%
%
%\ref{M1} and \ref{M2} promise one path of $\mathcal M^{*}$ per obligation, though, while a run of
%$\mathcal A_\varphi$ commits to one branch at every \ORn-node, and no single branch serves every
%obligation: under $\mathrm{AG}(\mathrm{EF}\,p\wedge\mathrm{EF}\,q)$ a node must head for $p$ on some
%visits and for $q$ on others. A location therefore pairs a saturated node with a \emph{schedule},
%which follows the finite path discharging one of the at most $n+1$ obligations pending there, the
%divergence of \ref{M1} and the until pursuits active at it, resolving the \ORn-nodes along that
%path and turning round-robin to the next obligation once the current one is discharged. Its memory
%is the pair (obligation, position on that path), so $S$ stays finite.
%

%
%Call a node of $\mathcal M^{*}$ a \emph{state} if it is
%saturated, so that its children are the successor steps; the remaining, \emph{expansion} edges
%neither elapse time nor change location, and reset a clock only through a \ruleref{launch} or a
%\ruleref{Ustopcarry}. At most $n+1$ obligations are pending at $N$, the divergence requirement of
%\ref{M1} and the until pursuits active there, and \ref{M1}, \ref{M2} give a finite path $\pi_{N,o}$
%of $\mathcal M^{*}$ discharging each. A \emph{schedule} selects at every \ORn-node the child
%$\pi_{N,o}$ prescribes and moves round-robin to the next pending obligation once $o$ is discharged,
%so its memory, the pair (obligation, position on $\pi_{N,o}$), is finite. The witness
%$\mathcal A_\varphi$ has as locations the pairs (state, schedule memory), labelled by
%$\Gamma_{\mathsf{prop}}$ and carrying $\Inv(N)$, and as edges the scheduled successor steps, guarded
%by the zone of the next state reached through expansion edges and resetting the clocks that segment
%resets.
%
\begin{lemma}[Recurrence]\label{lem:recurrence}
Let $\mu=\mathrm{A}(\varphi\,\mathbf{U}_\kI\,\psi)$ with $a=\inf\kI$, $b=\sup\kI<\infty$, and let
$\rho$ be a run of $\mathcal A_\varphi$ through a node carrying $\mu^{\ast}$, launched at $t_0$. Then
\ruleref{Ustopcarry} fires on $\rho$ at times $t_0<w_1<w_2<\cdots$ with $w_1-t_0\in\kI$ and
$w_{j+1}-w_j\le b-a$, and $\rho\models\varphi\,\mathbf{U}_\kI\,\psi$ from every position at or after
$t_0$.
\end{lemma}
\begin{proof}
No rule drops $\mu^{\ast}$: \ruleref{Ustopcarry} keeps it in $\Gamma$, \ruleref{Ugocarry} moves
$\AX\mu^{\ast}$ to $\Ctx$, whence every successor carries it again, and \ruleref{Absorb} bars a fresh
\ruleref{Cov} while the pursuit is live. So $\mu^{\ast}$ rides $\rho$ for ever, and with it the
deadline on $x_\mu$, which is $b$ until the first witness and $b-a$ after. Both are at most $M$, so
$x_\mu$ is bounded along $\rho$, and a cycle leaving a bounded clock unreset is not
time-progressing. By \ref{M1} the schedule reaches such a cycle and traverses it without end, and
while the pursuit is live only \ruleref{Ustopcarry} resets $x_\mu$, so the stop recurs: it fires at
times $t_0<w_1<w_2<\cdots$. At the first firing $x_\mu$ measures from $t_0$ and the guard is \kI, so
$w_1-t_0\in\kI$; at each later one it measures from the preceding witness and $\Deadline$ caps it at
$b-a$, so $w_{j+1}-w_j\le b-a$.

For the last claim let $t\ge t_0$ and let $w_j$ be the first witness at or after $t+a$. If $j=1$ then
$w_1\le t_0+b\le t+b$; otherwise $w_{j-1}<t+a$, so $w_j\le w_{j-1}+(b-a)<t+b$. Either way
$w_j-t\in\kI$, and $\psi$ holds at $w_j$ while \ruleref{Ugocarry} asserts $\varphi$ at every position
before it, which is \eqref{eq:sat_u}.
\end{proof}

\begin{theorem}\label{thm:soundness}
If Algorithm~\ref{alg:otf} returns \SAT{} then $\varphi$ is satisfiable.%, %and $\mathcal A_\varphi$ is a witness for it.
\end{theorem}
\begin{proof}
By Theorem~\ref{thm:alg} the root lies in $\mathcal M^{*}$, so $\mathcal A_\varphi$ is defined, and
Lemma~\ref{lem:zones}\ref{zoneitem:feasible} supplies valuations and delays along its paths. It is non-blocking: its only
invariants are the deadlines in $\Deadline$, and cap and stop guard derive from the same interval with the
same strictness, so the stop edge is enabled up to the cap and by \ref{M2} the schedule takes it. By
\ref{M1} the schedule reaches a time-progressing cycle, which Lemma~\ref{lem:zones}\ref{zoneitem:nonzeno} unfolds into
a run whose iterations each elapse a fixed $\delta>0$, so $\mathrm{Runs}^\infty$ is non-empty
everywhere, which is all the semantics asks of it.

We show $\mathcal T,(s_0,\mathbf 0)\models\varphi$ for the induced TTS by induction on $\varphi$
along the marking. Propositional goals hold by saturation and closure of contradictory nodes;
$\EX\psi$ obtains a successor satisfying $\psi$ from its \ruleref{EX} child, and $\AX\psi$ holds at
all successors since each child inherits $\Ctx$. For a timed modality demanded at $N$,
\ruleref{Cov} is conjunctive, so $\mathcal M^{*}$ retains both the discharge child, carrying the
remaining goals, and the launch, which establishes $\mu$ with $x_\mu$ reset, so \kI is measured
from this configuration and the pursuit rides every successor, as a goal under $\mathrm{E}$ and as
context under $\mathrm{A}$. A re-demand is covered anew, carried or absorbed, so every onset is established: by a pursuit
of its own, by the carried pursuit that answers all of them, or, when the window is unbounded, by
the witnesses that later launches are forced to produce. An existential until is witnessed by the path \ref{M2} supplies, extended past
its \ruleref{Ustop} by a divergent run; a universal one cannot be evaded, $\Inv(N)$ blocking time at
the cap while the pursuit is imposed on all successors, and when the demand is persistent this is
Lemma~\ref{lem:recurrence}. \ruleref{Rsplit} discharges a release by partitioning the zone
conjunctively: \ruleref{Runf} asserts $\psi$ in the window until $\varphi$ releases it, carries the
pursuit on before the window, and past it the obligation has lapsed.
Lemma~\ref{lem:zones}\ref{zoneitem:stability} lifts each conclusion to the positions between nodes.
\end{proof}

\begin{theorem}[Completeness]\label{thm:completeness}
If $\varphi$ is non-punctual and satisfiable then Algorithm~\ref{alg:otf} returns \SAT.
\end{theorem}
\begin{proof}
Completeness goes the other way: the realisable nodes of \S\ref{ss:nodes} are themselves a successful
marking. Writing $\mathcal R$ for them, we show $\mathcal R\subseteq F(\mathcal R)$; then $\mathcal R\subseteq\mathcal M^{*}$ by
Proposition~\ref{prop:gfp}, and $N_0\in \mathcal R$ as $\varphi$ is satisfiable, so Theorem~\ref{thm:alg}
returns \SAT.

\ref{D1}: a realisable node has a non-empty zone and a consistent $\Gamma_{\mathsf{prop}}$, and
\ruleref{Umiss} does not fire on it: a pending $\mu^\bullet$ is realisable at no valuation beyond
its window, so its zone meets $x_\mu\le b$.
\ref{D2}: let $N\in \mathcal R$ be realised by $\mathcal T$. All children of an \ANDn-rule are realised by the
same $\mathcal T$: \ruleref{Cov} because $\mathcal T$ satisfies $\Gamma\setminus\{\mu\}$ and $\mu$
alike, \ruleref{Rsplit} because its children partition $Z_N$ and realisability is monotone under
narrowing, \ruleref{Delay} and \ruleref{EX} because $\mathcal T$ admits the passage of time and
supplies a successor for every $\EX\psi$, each inheriting $\Ctx$. At an \ORn-rule $\mathcal T$ selects
a child: a disjunct it satisfies for \ruleref{Disj}; for \ruleref{Uunf} and \ruleref{Uunfcarry} the
stop if $\psi$ holds now with $x_\mu\in\kI$ and \ruleref{Ugo} otherwise, the witness being later and
$\varphi$ holding now; for \ruleref{Runf} the stop if $\varphi$ has
released the obligation and the recursive child otherwise. The zone operations being exact, no chosen
zone is empty. A demand re-imposed at a continuum of onsets is answered onset by onset, every re-demand being
covered anew, carried or absorbed according to whether the running pursuit reaches the new onset,
so the selections above serve each onset of each zone.
\ref{D3}: fix $N\in \mathcal R$ and a run $\rho$ of $\mathcal T$ from a configuration for $N$. It
is non-Zeno, so no clock stays both unreset and bounded along it, and following $\rho$ through the
choices above gives a path of $\mathcal R$ from $N$ which by Lemma~\ref{lem:finiteness} repeats a node; that
segment elapses positive time with every clock reset on it or extrapolated above $M$, hence is a
time-progressing cycle of $\mathcal R$ reached from $N$. An until pursuit active at $N$ is one $\mathcal T$
fulfils at a finite position of $\rho$, which the path following $\rho$ carries it to.

Non-punctuality enters at the repetition step: a node summarises the pending obligations of a
persistent universal until by $x_\mu$ alone, the time since the last witness capped at $b-a$, and by
Lemma~\ref{lem:recurrence} that summary is faithful exactly when witnesses spaced at most $b-a$ apart
answer every onset, which needs $b>a$.
\end{proof}


\begin{corollary}\label{cor:depthone}
Satisfiability of non-punctual TCTL is \textsc{ExpTime}-complete, and the procedure is sound and
complete for TCTL of temporal nesting depth at most one, punctual intervals included.
\end{corollary}
\begin{proof}
Membership is Theorems~\ref{thm:alg}--\ref{thm:completeness}; hardness is inherited from CTL
satisfiability~\cite{Clarke99} by fixing every interval to $[0,\infty)$. At depth one every argument
is propositional, so no modality is demanded inside the pursuit of another and no fresh onset arises;
the obligations live at a node are exactly the pursuits it carries, recorded in $\Gamma$, $\Ctx$ and
$\Deadline$ with their elapsed time in $Z$, so a repeat agrees on the full obligation state and
Theorem~\ref{thm:completeness} goes through for every \kI.
\end{proof}

\noindent
Depth one covers reachability, safety and inevitability, the queries
\textsc{Uppaal}~\cite{BehrmannLarsenYi2004Uppaal} and \textsc{Romeo}~\cite{lime-TACAS-09} are asked
for; the one common query beyond it, leads-to $\mathrm{AG}(p\rightarrow\mathrm{AF}_\kI q)$, stays
covered by Theorem~\ref{thm:completeness} whenever \kI is non-punctual.
